\documentclass[11pt]{article}
\usepackage[a4paper,margin=28mm]{geometry}
\usepackage{amsmath,amssymb,amsthm}
\usepackage[colorlinks=true,linkcolor=blue,urlcolor=blue,citecolor=blue]{hyperref}
\newtheorem*{theorem*}{Theorem}
\newtheorem*{corollary*}{Corollary}
\newenvironment{statement}[1]{\par\medskip\noindent\phantomsection\textbf{#1}\ }{\par\medskip}
\setlength{\emergencystretch}{3em}
\newcommand{\Q}{\mathbb Q}
\newcommand{\Z}{\mathbb Z}
\newcommand{\F}{\mathbb F}
\DeclareMathOperator{\End}{End}
\DeclareMathOperator{\Gal}{Gal}
\DeclareMathOperator{\GL}{GL}
\DeclareMathOperator{\PGL}{PGL}
\DeclareMathOperator{\Hom}{Hom}
\DeclareMathOperator{\NS}{NS}
\DeclareMathOperator{\Sym}{Sym}
\DeclareMathOperator{\tr}{tr}
\title{No rational \(17\)-torsion on geometrically split abelian surfaces}
\author{}
\date{10 November 2026}

\begin{document}
\maketitle

\begin{theorem*}
Let \(A/\Q\) be a geometrically split abelian surface. Then \(A(\Q)\) contains no point of
order \(17\).
\end{theorem*}

\begin{corollary*}
No geometrically split genus-\(2\) Jacobian over \(\Q\) has a rational point of order \(17\).
Equivalently, a genus-\(2\) Jacobian over \(\Q\) with a rational point of order \(17\) is
geometrically simple. In particular,
\href{https://genus-2-torsion.github.io/group.html?g=17}{\([17]\)} does not occur as the rational
torsion subgroup of a geometrically split genus-\(2\) Jacobian over \(\Q\).
\end{corollary*}

Here ``geometrically split'' means that \(A_{\bar\Q}\) is isogenous to a product of two elliptic
curves. We impose no restriction on the degree of this isogeny, on the field of definition of the
geometric endomorphisms, or on the integral endomorphism ring; we do not assume that \(A\) is
principally polarized, and we do not assume that \(A\) is simple over \(\Q\). For surfaces that
are isogenous over \(\Q\) to a product of two elliptic curves the theorem follows from Mazur's
theorem
(\href{https://genus-2-torsion.github.io/proofs/qsplit-prime-obstruction.html}{prime torsion
obstruction}); the new case is that of surfaces that are simple over \(\Q\). The theorem contains
the statement about order \(34\) of the
\href{https://genus-2-torsion.github.io/proofs/geometric-33-34-obstruction.html}{proof for orders
33 and 34}.

The proof distinguishes cases according to the geometric endomorphism algebra
(Section~\ref{cases}). Most cases are settled by the known results on torsion of elliptic curves
over fields of degree at most \(2\) (Section~\ref{small-degree}), by counting points on a
reduction at \(2\) (Section~\ref{good-reduction}), or by comparing lattices in the \(17\)-adic
Tate module (Section~\ref{lattices}). In the one remaining case the geometric endomorphisms are
defined over a cyclic quartic field. There a rational point of order \(17\) produces a
non-cuspidal point of the modular curve \(X_0(34)\) over a \emph{real} quadratic field \(K\)
which the Atkin--Lehner involution \(w_2\) maps to its conjugate (Section~\ref{c4}). The
rational points of \(X_0(34)/w_2\) show that such a point exists only for
\(K=\Q(\sqrt{-15})\) (Section~\ref{x034}).

The Magma script and its log cited below are published with this proof, as
\href{https://genus-2-torsion.github.io/data/knowledge/sources/geometric_17_checks.m}{\texttt{geometric\_\allowbreak{}17\_\allowbreak{}checks.m}}
and
\href{https://genus-2-torsion.github.io/data/knowledge/sources/geometric_17_checks.log}{\texttt{geometric\_\allowbreak{}17\_\allowbreak{}checks.log}}.

\section{Conventions and standard facts}\label{facts}

All number fields are subfields of a fixed algebraic closure \(\bar\Q\) of \(\Q\). For a field
\(k\) we write \(G_k=\Gal(\bar k/k)\). For a prime \(\ell\), \(\chi_\ell\) is the \(\ell\)-adic
cyclotomic character and \(\bar\chi_\ell\) its reduction modulo \(\ell\); thus
\(\bar\chi_{17}:G_{\Q}\to\F_{17}^\times\) is surjective of order \(16\). For an abelian variety
\(X\) over a field \(k\) of characteristic \(0\), \(T_\ell X\) is its \(\ell\)-adic Tate module,
\(V_\ell X=T_\ell X\otimes\Q_\ell\), \(\End^0=\End\otimes\Q\), \(\Hom^0=\Hom\otimes\Q\), and
\(X\sim_k Y\) means that \(X\) and \(Y\) are isogenous over \(k\). A quasi-isogeny is an
invertible element of \(\Hom^0\). An elliptic curve is ``without CM'' if its endomorphism ring
over \(\bar\Q\) is \(\Z\). For a \(G\)-stable \(\Z_\ell\)-lattice \(\Lambda\) in a
\(\Q_\ell[G]\)-module we put \(\bar\Lambda=\Lambda/\ell\Lambda\), and \((\cdot)^{\mathrm{ss}}\)
denotes semisimplification. A ``constituent'' is a Jordan--H\"older factor.

\emph{Conjugates.} For a variety \(Y\) over \(\bar\Q\) and \(g\in G_{\Q}\), \(Y^g\) is the base
change of \(Y\) along \(g:\bar\Q\to\bar\Q\), and \(g:Y(\bar\Q)\to Y^g(\bar\Q)\) is the induced
bijection. A morphism \(f:Y\to Z\) has the conjugate \(g(f):Y^g\to Z^g\), characterized by
\(g(f)(g(P))=g(f(P))\). We have \((Y^g)^h=Y^{hg}\). If \(Y\) is defined over a number field
\(k\), then \(Y^g\) depends only on the restriction of \(g\) to \(k\), and \(Y^g=Y\) for
\(g\in G_k\); a morphism \(f\) between varieties defined over \(k\) is defined over \(k\) if and
only if \(g(f)=f\) for all \(g\in G_k\). Conjugation is compatible with composition, with duals
of isogenies and with degrees.

For a character \(\eta\) of \(G_K\), where \(K\) is a number field, and \(\tau\in G_{\Q}\)
normalizing \(G_K\), we put \(\eta^\tau(g)=\eta(\tau^{-1}g\tau)\); similarly, for a
\(G_K\)-module \(M\), \(M^\tau\) is \(M\) with \(g\) acting as \(\tau^{-1}g\tau\). If \(X\) is an
abelian variety over \(K\), then \(\tau\) induces an isomorphism of \(G_K\)-modules
\(X[n]^\tau\cong X^\tau[n]\), because \(\tau(g'v)=(\tau g'\tau^{-1})\tau(v)\) for \(g'\in G_K\).

We use the following standard facts.
\begin{itemize}
\item \textbf{(F1) Degrees.} The degree extends multiplicatively to quasi-isogenies between
abelian varieties of dimension \(d\), by \(\deg(f/n)=\deg(f)/n^{2d}\). Let \(X\) be an elliptic
curve over \(\bar\Q\) without CM, so that \(\End(X^2)=M_2(\Z)\). An element \(N\in M_2(\Z)\) with
\(\det N\neq0\) has degree \((\det N)^2\): write \(N=U\,\mathrm{diag}(d_1,d_2)\,U'\) with
\(U,U'\in\GL_2(\Z)\); the outer factors are automorphisms and the middle one has degree
\(d_1^2d_2^2\). Hence \(\deg N=(\det N)^2\) for all \(N\in\GL_2(\Q)\). If \(\mu:X'\to X\) is an
isogeny of elliptic curves, then \(\mathrm{diag}(\mu,\mu):X'^2\to X^2\) has degree
\((\deg\mu)^2\).

\item \textbf{(F2) Brauer--Nesbitt.} For a continuous representation of a profinite group \(G\)
on a finite-dimensional \(\Q_\ell\)-vector space \(V\), the semisimplification
\(\bar\Lambda^{\mathrm{ss}}\) does not depend on the \(G\)-stable lattice \(\Lambda\subset V\).
The same applies to two isomorphic representations, after transporting a lattice through the
isomorphism. Indeed, after scaling we may assume \(\Lambda'\subseteq\Lambda\); the \(G\)-stable lattices
\(M_i=\Lambda'+\ell^i\Lambda\) satisfy \(M_0=\Lambda\), \(M_i=\Lambda'\) for large \(i\), and
\(\ell M_i\subseteq M_{i+1}\subseteq M_i\). For \(G\)-stable lattices \(M'\subseteq M\) with
\(\ell M\subseteq M'\), both \(M/\ell M\) and \(M'/\ell M'\) have the composition factors of
\(M/M'\) and of \(M'/\ell M\): the first contains \(M'/\ell M\) with quotient \(M/M'\), the second
contains \(\ell M/\ell M'\cong M/M'\) with quotient \(M'/\ell M\).

\item \textbf{(F3) Forms of a square.} Let \(k\) be a number field, \(E\) an elliptic curve over
\(k\) without CM, \(B\) an abelian surface over \(k\), and \(f:E^2_{\bar k}\to B_{\bar k}\) a
quasi-isogeny. For \(g\in G_k\) put
\[\rho(g)=f^{-1}\circ g(f)\in\operatorname{Aut}^0(E^2_{\bar k})=\GL_2(\Q).\]
\begin{enumerate}
\item[(i)] \(\rho:G_k\to\GL_2(\Q)\) is a homomorphism with finite image. Indeed
\(\rho(gh)=f^{-1}g(f)\cdot g\bigl(f^{-1}h(f)\bigr)=\rho(g)\rho(h)\), because the elements of
\(M_2(\Q)=\End^0(E^2_{\bar k})\) are defined over \(k\); and \(\rho\) is trivial on \(G_{k'}\)
for a finite extension \(k'\) of \(k\) over which \(f\) is defined.
\item[(ii)] For \(\varphi\in\End^0(B_{\bar k})\) we have
\(f^{-1}g(\varphi)f=\rho(g)\,(f^{-1}\varphi f)\,\rho(g)^{-1}\), as one sees by applying \(g\) to
\(f^{-1}\varphi f\in M_2(\Q)\). So, in the identification \(\End^0(B_{\bar k})\cong M_2(\Q)\)
given by \(f\), the element \(g\) acts by \(\mathrm{Ad}(\rho(g))\), and \(\End^0_k(B)\)
corresponds to the commutant of \(\rho(G_k)\) in \(M_2(\Q)\).
\item[(iii)] \(f\) induces an isomorphism \(V_\ell(E)\otimes_{\Q}\Q^2=V_\ell(E^2)\to V_\ell(B)\)
under which \(g\in G_k\) acts on the left-hand side by \(\rho_{E,\ell}(g)\otimes\rho(g)\):
indeed \(g(f(v))=g(f)(gv)=f\bigl(\rho(g)(gv)\bigr)\).
\end{enumerate}

\item \textbf{(F4) Elliptic curves over fields of degree at most \(2\).} No elliptic curve over
\(\Q\) or over a quadratic field has a rational point of order \(17\). Over \(\Q\) this is part
of Mazur's theorem \cite[Theorem~(8)]{Maz}. Over quadratic fields it is Kamienny's theorem
\cite[Theorem~3.1]{Kam}: the order of the torsion subgroup of an elliptic curve over a quadratic field is not
divisible by any prime greater than \(13\). (In the notation of \cite[Section~1]{DKSS}, where
\(S(d)\) is the set of primes that occur as the order of a point on an elliptic curve over a
number field of degree \(d\): \(S(1)=\mathrm{Primes}(7)\) and \(S(2)=\mathrm{Primes}(13)\).)
\end{itemize}

\begin{statement}{Lemma 1.}\label{lemma-sequence}
Let \(X\) be an abelian surface over a number field \(k\), let \(E\subset X\) be an elliptic
curve defined over \(k\), and let \(Q\in X(k)\) have prime order \(q\). Then either
\(Q\in E(k)\), or the image of \(Q\) in \((X/E)(k)\) has order \(q\).
\end{statement}

\emph{Proof.} The image of \(Q\) is zero exactly when \(Q\in E\). \qed

\section{The case distinction}\label{cases}

Let \(A/\Q\) be a geometrically split abelian surface, with \(A_{\bar\Q}\sim E_1\times E_2\).
Exactly one of the following holds:
\begin{itemize}
\item \textbf{(N)} \(E_1\) and \(E_2\) are not isogenous;
\item \textbf{(CM)} \(E_1\sim E_2=:E\) and \(E\) has complex multiplication;
\item \textbf{(S)} \(E_1\sim E_2=:E\) and \(E\) has no complex multiplication.
\end{itemize}

In case (S), \(D:=\End^0(A_{\bar\Q})\cong M_2(\Q)\). Let \(L\) be the smallest field over which
all geometric endomorphisms of \(A\) are defined; it is a finite Galois extension of \(\Q\). The
group \(\Gamma=\Gal(L/\Q)\) acts faithfully on \(D\), and \(\End^0_{\Q}(A)=D^\Gamma\). By
Skolem--Noether this action embeds \(\Gamma\) into \(\operatorname{Aut}(M_2(\Q))=\PGL_2(\Q)\);
we write \(g\mapsto[R_g]\), so that \(g\) acts on \(D\) by \(\mathrm{Ad}(R_g)\); for
\(g\in G_{\Q}\) we put \(R_g=R_{g|_L}\). The trace and
the determinant of an element of \(D\) are those of the corresponding matrix; they do not depend
on the isomorphism \(D\cong M_2(\Q)\).

A finite subgroup of \(\PGL_2(\mathbb R)\) is cyclic or dihedral: \(\PGL_2(\mathbb R)\) acts on the
hyperbolic plane by isometries (an element of negative determinant by
\(z\mapsto(a\bar z+b)/(c\bar z+d)\)), a finite subgroup fixes a point, and the stabilizer of a
point is isomorphic to \(O(2)\). If \([R]\in\PGL_2(\Q)\) has order \(n>2\), the ratio
of the eigenvalues of \(R\) is a primitive \(n\)-th root of unity \(\zeta\), and
\((\tr R)^2/\det R=2+\zeta+\zeta^{-1}\in\Q\); hence \(n\in\{3,4,6\}\). Therefore
\[\Gamma\in\{C_1,\ C_2,\ C_3,\ C_4,\ C_6,\ V_4,\ D_6,\ D_8,\ D_{12}\},\]
where \(D_{2n}\) is dihedral of order \(2n\) (so \(D_6\cong S_3\)) and \(V_4\) is the Klein four
group. An element \(r\) of order \(n\in\{3,4,6\}\) has a unique representative \(R\) with
\(\tr R=1\) (the trace is nonzero, since otherwise \(R^2\) would be a scalar); its
characteristic polynomial is \(t^2-t+1/c\) with \(c=1,2,3\) respectively, of discriminant
\(-3,-1,-1/3\). Since \(R(1-R)=1/c\), the trace-one representative of \(r^{-1}\) is \(1-R\).

Every elliptic subvariety of \(A_{\bar\Q}\) is the image of an integral multiple of a rank-one
idempotent of \(D\), so it is defined over \(L\). If \(E,E'\subset A_L\) are elliptic
subvarieties, every geometric homomorphism \(E'\to E\) is defined over \(L\): composing with an
\(L\)-rational quasi-projection \(A\to E'\) and the inclusion \(E\subset A\) embeds
\(\Hom^0(E'_{\bar\Q},E_{\bar\Q})\) \(G_L\)-equivariantly into \(D\), on which \(G_L\) acts
trivially.

The Theorem is proved case by case: type (N), and type (S) with \(\Gamma\in\{C_1,C_2\}\), in
Section~\ref{small-degree}; type (CM) in Section~\ref{good-reduction}; type (S) with
\(\Gamma\in\{C_3,C_6,V_4,D_6,D_8,D_{12}\}\) in Section~\ref{lattices}; and type (S) with
\(\Gamma=C_4\) in Sections~\ref{c4} and~\ref{x034}. The proof is assembled in
Section~\ref{census}.

\section{Elliptic curves over fields of degree at most \(2\)}\label{small-degree}

\begin{statement}{Lemma 2.}\label{lemma-small}
Let \(A\) be an abelian surface over \(\Q\) and \(k\) a number field with \([k:\Q]\le2\) such
that \(A_k\) contains an elliptic curve defined over \(k\). Then \(A(\Q)\) contains no point of
order \(17\).
\end{statement}

\emph{Proof.} Let \(E\subset A_k\) be an elliptic curve defined over \(k\), and suppose that
\(Q\in A(\Q)\) has order \(17\). By \hyperref[lemma-sequence]{Lemma~1} over \(k\), either
\(Q\in E(k)\) or the image of \(Q\) in \((A_k/E)(k)\) has order \(17\). Both \(E\) and \(A_k/E\)
are elliptic curves over a field of degree at most \(2\), which contradicts (F4). \qed

\medskip
\textbf{Type (N).} Let \(A\) be of type (N). Then \(A_{\bar\Q}\) has exactly two elliptic
subvarieties: an elliptic subvariety is isogenous to \(E_1\) or to \(E_2\), and for each \(i\)
the space \(\Hom^0(E_i,A_{\bar\Q})\) has rank one over the field \(\End^0(E_i)\), so all nonzero
maps \(E_i\to A_{\bar\Q}\) have the same image \(B_i\); and \(B_1\neq B_2\) because they are not
isogenous. The group \(G_{\Q}\) permutes \(B_1\) and \(B_2\). The stabilizer of \(B_1\) is an
open subgroup of index at most \(2\); its fixed field \(k\) has degree at most \(2\), and \(B_1\)
is defined over \(k\). By \hyperref[lemma-small]{Lemma~2}, \(A(\Q)\) has no point of order
\(17\).

\medskip
\textbf{Type (S) with \(\Gamma\in\{C_1,C_2\}\).} Here \([L:\Q]\le2\), and \(A_{\bar\Q}\sim E^2\)
contains an elliptic subvariety, which is defined over \(L\) by Section~\ref{cases}. By
\hyperref[lemma-small]{Lemma~2} with \(k=L\), \(A(\Q)\) has no point of order \(17\).

\medskip
This argument does not extend to all the remaining groups \(\Gamma\): for \([L:\Q]=4\) the
analogue of (F4) fails, since elliptic curves over quartic fields can have a point of order
\(17\) (see Remark~1 in Section~\ref{remarks}).

\section{Potentially good reduction at \(2\)}\label{good-reduction}

\begin{statement}{Proposition 3.}\label{prop-good}
Let \(A\) be an abelian surface over \(\Q\) with potentially good reduction at \(2\). Then
\(A(\Q)\) contains no point of order \(17\).
\end{statement}

Here ``potentially good reduction at \(2\)'' means that \(A\) acquires good reduction over a
finite extension \(F'\) of \(\Q_2\); for example \(F'=k_v\) when \(A_k\) has good reduction at a
prime \(v\mid2\) of a number field \(k\). We fix an embedding \(\bar\Q\subset\bar\Q_2\), so that
\(G_{\Q_2}\subseteq G_{\Q}\).

\emph{Proof.}
\begin{enumerate}
\item \emph{Good reduction over a totally ramified extension of \(\Q_2\).} Let
\(I\subset G_{\Q_2}\) be the inertia group. By the N\'eron--Ogg--Shafarevich criterion
\cite[Theorem~1]{ST}, the inertia group \(I\cap G_{F'}\) of \(F'\) acts trivially on \(T_3A\); as it
has finite index in \(I\), the group \(I\) acts on \(T_3A\) through a finite quotient. Put \(N=I\cap\ker(\rho_{A,3})\), an open normal subgroup of \(I\) that is normal
in \(G_{\Q_2}\). The group \(G_{\Q_2}/N\) is an extension of \(G_{\Q_2}/I\cong\hat\Z\) by the
finite group \(I/N\). Since \(\hat\Z\) is a free profinite group on one generator, the extension
splits: mapping the generator to a lift of the Frobenius element defines a continuous section.
So there is a closed subgroup \(H\supseteq N\) of \(G_{\Q_2}\) with \(H\cap I=N\) and
\(HI=G_{\Q_2}\); it has index \([I:N]\). Its fixed field \(F\) is a finite extension of \(\Q_2\)
with residue field \(\F_2\), i.e. totally ramified, and its inertia group \(N\) acts trivially
on \(T_3A\). By the same criterion, \(A\) has
good reduction \(\bar A/\F_2\) over \(F\). Prime-to-\(2\) torsion of \(A(F)\) injects into
\(\bar A(\F_2)\): for odd \(m\), the \(m\)-torsion of the N\'eron model of \(A_F\), an abelian
scheme over the ring of integers \(\mathcal O_F\), is finite \'etale over \(\mathcal O_F\), so
reduction is injective on \(A(F)[m]\). So, if \(A(\Q)\) has a point of order \(17\), then \(17\) divides
\(\#\bar A(\F_2)=P(1)\), where \(P\) is the characteristic polynomial of the Frobenius
endomorphism \(\pi\) of \(\bar A\).

\item \emph{The Weil polynomial.} The polynomial \(P\in\Z[T]\) is monic of degree \(4\), its
complex roots have absolute value \(\sqrt2\) by Weil's theorem, and \(P(0)=\deg\pi=4\). The
non-real roots form pairs \(\omega,\bar\omega\) with \(\omega\bar\omega=2\); each pair
contributes a factor \(T^2-\beta T+2\) with \(\beta=\omega+\bar\omega\) real and
\(|\beta|<2\sqrt2\). The possible real roots \(\sqrt2\) and \(-\sqrt2\) are conjugate over
\(\Q\), so they have the same multiplicity \(m\in\{0,1,2\}\). If \(m=1\), then
\(P=(T^2-2)(T^2-\beta T+2)\) and \(P(0)=-4\), which is excluded. If \(m=2\), then
\(P=(T^2-2)^2=(T^2-2\sqrt2\,T+2)(T^2+2\sqrt2\,T+2)\). In all cases
\[P(T)=(T^2-\beta_1T+2)(T^2-\beta_2T+2)=T^4+aT^3+bT^2+2aT+4\]
with real numbers \(|\beta_i|\le2\sqrt2\) and integers \(a=-(\beta_1+\beta_2)\),
\(b=\beta_1\beta_2+4\), and
\[P(1)=5+3a+b=(3-\beta_1)(3-\beta_2).\]

\item \emph{\(P(1)\) is not divisible by \(17\).} Put \(x_i=3-\beta_i\). Then
\(0<3-2\sqrt2\le x_i\le3+2\sqrt2\), so \(P(1)=x_1x_2\) is a positive integer and
\(P(1)\le(3+2\sqrt2)^2=17+12\sqrt2<34\), because \(288<289\). Suppose that \(17\) divides
\(P(1)\). Then \(x_1x_2=17\), and both \(x_i\) lie in the interval
\([\,17/(3+2\sqrt2),\ 3+2\sqrt2\,]=[\,51-34\sqrt2,\ 3+2\sqrt2\,]\). The function
\(x\mapsto x+17/x\) is convex on this interval and takes the same value \(54-32\sqrt2\) at both
endpoints, which are exchanged by \(x\mapsto17/x\); and it is at least \(2\sqrt{17}\). Hence the
integer \(x_1+x_2=6+a\) satisfies
\[8<2\sqrt{17}\ \le\ 6+a\ \le\ 54-32\sqrt2<9,\]
where the outer inequalities are \(64<68\) and \(2025<2048\). This is impossible. \qed
\end{enumerate}

As a check, part (A) of the script
\href{https://genus-2-torsion.github.io/data/knowledge/sources/geometric_17_checks.m}{\texttt{geometric\_\allowbreak{}17\_\allowbreak{}checks.m}}
(see Section~\ref{x034}) enumerates all pairs of integers \((a,b)\) for which
\(T^4+aT^3+bT^2+2aT+4\) has all its roots of absolute value \(\sqrt2\). There are \(35\) of them,
and the values of \(P(1)\) are \(1,\dots,16,19,20,25\).

\medskip
\textbf{Type (CM).} Let \(A\) be of type (CM). The curve \(E\) has potentially good reduction
everywhere, because it has CM (its \(j\)-invariant is an algebraic integer; see \cite[\S2]{ST}). Over a suitable
number field \(k\), \(A_k\) is isogenous to \(E_k^2\) and \(E_k\) has good reduction at the
primes above \(2\). For a prime \(\ell\neq2\), \(V_\ell(A_k)\cong V_\ell(E_k)^2\) as
\(G_k\)-modules, and this is unramified at every prime \(v\mid2\) of \(k\); so \(A_k\) has good
reduction at \(v\) by the N\'eron--Ogg--Shafarevich criterion \cite[Theorem~1]{ST}. So \(A\) has
potentially good reduction at \(2\), and by \hyperref[prop-good]{Proposition~3}, \(A(\Q)\) has no point of order
\(17\).

\section{Type (S): lattices in the Tate module}\label{lattices}

In this section \(A\) is of type (S), \(p\) is an odd prime, \(V=V_p(A)\) and \(T=T_p(A)\). The
group \(G_{\Q}\) acts on \(D\otimes\Q_p\) through \(\Gamma\), and \(g(xv)=g(x)\,g(v)\) for
\(x\in D\otimes\Q_p\) and \(v\in V\). We fix an isomorphism
\(D\otimes\Q_p\cong\End_{\Q_p}(S)\) with \(S=\Q_p^2\); as a \(D\otimes\Q_p\)-module, \(V\) is
isomorphic to \(S\oplus S\).

\begin{statement}{Lemma 4 (polarization).}\label{lemma-pol}
Let \(T'\subset V\) be a \(G_{\Q}\)-stable lattice such that \((T'/pT')|_{G_L}\) has trivial
semisimplification. Then \(\Q(\zeta_p)\subseteq L\).
\end{statement}

\emph{Proof.} A polarization of \(A\) defined over \(\Q\) gives an isomorphism of
\(G_{\Q}\)-modules \(V\cong\Hom(V,\Q_p(1))\). Transport the lattice \(\Hom(T',\Z_p(1))\) into
\(V\). Its reduction is \(\Hom(T'/pT',\mu_p)\), whose semisimplification as a \(G_L\)-module
consists of four copies of \(\bar\chi_p|_{G_L}\), while that of \(T'/pT'\) consists of four
copies of the trivial character. By (F2) over \(G_L\), \(\bar\chi_p|_{G_L}=1\), that is, \(\Q(\zeta_p)\subseteq L\), since \(\Q(\zeta_p)\) is the
fixed field of \(\ker\bar\chi_p\). The degree of the
polarization plays no role. \qed

\begin{statement}{Lemma 5 (unramified order).}\label{lemma-unram}
Suppose that \(\Gamma\) normalizes an unramified quadratic \(\Z_p\)-order
\(\mathcal B\subset D\otimes\Q_p\) (the ring of integers of an unramified quadratic extension of \(\Q_p\)
contained in \(D\otimes\Q_p\)), that some \(g_0\in\Gamma\) acts on \(D\otimes\Q_p\) by
\(\mathrm{Ad}(R_0)\) for a unit \(R_0\in\mathcal B^\times\) whose reduction lies in
\(\F_{p^2}\setminus\F_p\), and that \(A(\Q)[p]\neq0\). Then \(\Q(\zeta_p)\subseteq L\).
\end{statement}

\emph{Proof.}
\begin{enumerate}
\item \emph{A \(\Gamma\)-stable maximal order.} The field \(\mathcal B\otimes\Q_p\) acts on \(S\), which
is a one-dimensional vector space over it. Since \(\mathcal B\) is a discrete valuation ring with
uniformizer \(p\), the \(\mathcal B\)-stable lattices in \(S\) are the lattices \(p^k\Lambda\),
\(k\in\Z\), where \(\Lambda=\mathcal Bv_0\) for a fixed nonzero \(v_0\in S\). For \(g\in\Gamma\) choose
any \(R_g\in(D\otimes\Q_p)^\times\) inducing the action of \(g\). Then \(R_g\Lambda\) is stable under
\(R_g\mathcal BR_g^{-1}=\mathcal B\), so \(R_g\Lambda=p^{k(g)}\Lambda\), and \(M_g:=p^{-k(g)}R_g\) lies in
\(\GL(\Lambda)\). It induces the action of \(g\) on \(D\otimes\Q_p\) and is determined by \(g\)
up to a factor in \(\Z_p^\times\). Consequently \(\Gamma\) preserves the maximal order
\(\mathcal R=\End(\Lambda)\cong M_2(\Z_p)\), which contains \(\mathcal B\), and
\(g\mapsto[\bar M_g]\in\PGL(U)\) is a homomorphism, where \(U=\Lambda/p\Lambda\) and
\(\bar M_g\) is the reduction of \(M_g\).

\item \emph{Tensor decomposition.} Put \(T'=\mathcal RT\). It is a \(G_{\Q}\)-stable lattice in
\(V\) and an \(\mathcal R\)-module. Choose a basis \(\lambda_1,\lambda_2\) of \(\Lambda\) and let
\(e_{ik}\in\mathcal R\) be the matrix units, \(e_{ik}\lambda_l=\delta_{kl}\lambda_i\). Then
\(T'=e_{11}T'\oplus e_{22}T'\), and \(e_{21}:e_{11}T'\to e_{22}T'\) is an isomorphism with
inverse \(e_{12}\). So \(W_0=e_{11}T'\) is free of rank \(2\) over \(\Z_p\), and
\[\Lambda\otimes_{\Z_p}W_0\to T',\qquad\lambda_l\otimes w\mapsto e_{l1}w,\]
is an isomorphism of \(\mathcal R\)-modules. Let \(\bar V=T'/pT'\) and \(W=W_0/pW_0\). Then
\(\bar V\cong U\otimes W\) as modules over \(\mathcal R/p\mathcal R=\End(U)\). For
\(g\in G_{\Q}\), let \(\bar g\) be its action on \(\bar V\), and write \(M_g\) for the matrix
attached to the image of \(g\) in \(\Gamma\). From \(g(xv)=M_gxM_g^{-1}\,g(v)\) for
\(x\in\mathcal R\) and \(v\in T'\) we get
\(\bar g\circ(\bar x\otimes1)=(\bar M_g\bar x\bar M_g^{-1}\otimes1)\circ\bar g\) for all
\(\bar x\in\End(U)\). Hence \((\bar M_g^{-1}\otimes1)\circ\bar g\) commutes with
\(\End(U)\otimes1\), so it is of the form \(1\otimes N_g\) with \(N_g\in\GL(W)\), and
\(\bar g=\bar M_g\otimes N_g\). Let
\[\tilde G=\{(g,M)\in G_{\Q}\times\GL(U)\ :\ M\in\F_p^\times\bar M_g\}.\]
This is a group, by step~1. For \((g,M)\in\tilde G\) there is a unique \(N\in\GL(W)\) with
\(\bar g=M\otimes N\), and \((g,M)\mapsto N\) is a homomorphism, since
\(\bar g\bar g'=MM'\otimes NN'\). So \(U\) and \(W\) are two-dimensional \(\tilde G\)-modules,
and \(\bar V\cong U\otimes W\) as \(\tilde G\)-modules, where \(\tilde G\) acts on \(\bar V\)
through its projection to \(G_{\Q}\).

\item \emph{\(U\) is irreducible.} We may take \(M_{g_0}=R_0\), since \(R_0\Lambda=\Lambda\). So
\(\tilde G\) contains an element acting on \(U\cong\mathcal B/p\mathcal B=\F_{p^2}\) as multiplication by an
element of \(\F_{p^2}\setminus\F_p\), which has no eigenvector in \(U\).

\item \emph{\(W\) is irreducible.} The lattices \(T\) and \(T'\) have the same reduction up to
semisimplification, by (F2). As \(A(\Q)[p]\neq0\), the trivial character is a constituent of
\(T/pT\), hence of \(\bar V\). If \(W\) had a \(\tilde G\)-stable line \(W_1\), then
\(U\otimes W_1\) and \(U\otimes W/W_1\) would be irreducible \(\tilde G\)-modules of dimension \(2\),
and \(\bar V\) would have no one-dimensional \(\tilde G\)-constituent, hence no one-dimensional
\(G_{\Q}\)-constituent: \(\tilde G\) maps onto \(G_{\Q}\), since it contains \((g,\bar M_g)\) for
every \(g\), and it acts on \(\bar V\) through this map.

\item \emph{\(G_L\) acts trivially on \(\bar V\).} For \(h\in G_L\), \(\bar M_h\) is a scalar, so
\(H=\{(h,1):h\in G_L\}\) is a subgroup of \(\tilde G\). It is normal, because \(G_L\) is normal
in \(G_{\Q}\), and it acts trivially on \(U\). By Clifford's theorem, \(W|_H\) is semisimple: the
\(\tilde G\)-translates of a simple \(H\)-submodule are simple \(H\)-submodules whose sum is a
nonzero \(\tilde G\)-submodule, hence all of \(W\). Since
\(\bar V|_H\cong W|_H\oplus W|_H\) has a trivial constituent (the restriction of a trivial
constituent of \(\bar V\)) and \(W|_H\) is semisimple, \(W^H\neq0\); as \(W^H\) is
\(\tilde G\)-stable and \(W\) is irreducible, \(W^H=W\). So \(G_L\) acts trivially on
\(\bar V=T'/pT'\).

\item \hyperref[lemma-pol]{Lemma~4} now gives \(\Q(\zeta_p)\subseteq L\). \qed
\end{enumerate}

\begin{statement}{Lemma 6 (split blocks).}\label{lemma-split}
Suppose that \(\Gamma\) normalizes a split \'etale order
\(\mathcal B=\Z_pe_1\oplus\Z_pe_2\subset D\otimes\Q_p\), where \(e_1,e_2\) are nonzero orthogonal
idempotents with \(e_1+e_2=1\), and that some element of \(\Gamma\) exchanges \(e_1\) and
\(e_2\). Let \(K\subset L\) be the fixed field of the subgroup of index \(2\) of \(\Gamma\) that
fixes \(e_1\). Then \(\Gal(L/K)\) acts on the line \(e_1(D\otimes\Q_p)e_2\) by a character
\(\eta\). If \(\eta\neq1\) and \(A(\Q)[p]\neq0\), then \(\Q(\zeta_p)\subseteq L\).
\end{statement}

\emph{Proof.} Every element of \(\Gamma\) permutes the two primitive idempotents \(e_1,e_2\) of
\(\mathcal B\); the elements fixing \(e_1\) also fix \(e_2=1-e_1\) and hence preserve the line
\(e_1(D\otimes\Q_p)e_2\). The character \(\eta\) has finite order, so its values are roots of
unity in \(\Z_p^\times\), on which reduction modulo \(p\) is injective; thus
\(\bar\eta\neq1\). We regard \(\eta\) as a character of \(G_K\).

Put \(T'=\mathcal BT=T_1\oplus T_2\) with \(T_i=e_iT\); it is a \(G_{\Q}\)-stable lattice. Let
\(M_i=T_i/pT_i\) and \(V_i=T_i\otimes\Q_p=e_iV\), of dimension \(2\). The group \(G_K\) preserves
each \(T_i\), and an element \(\tau\in G_{\Q}\setminus G_K\) exchanges them, so that
\(M_2\cong M_1^{\tau}\). A nonzero \(\xi\in e_1(D\otimes\Q_p)e_2\) maps \(V_2\) isomorphically
onto \(V_1\) and satisfies \(g(\xi v)=\eta(g)\,\xi\,g(v)\) for \(g\in G_K\). So \(\xi(T_2)\) is a
\(G_K\)-stable lattice in \(V_1\) with reduction \(\bar\eta\otimes M_2\), and by (F2),
\(M_1^{\mathrm{ss}}\cong\bar\eta\otimes M_2^{\mathrm{ss}}\).

The rational point of order \(p\) gives a trivial constituent of \(T/pT\), hence, by (F2), a
trivial \(G_K\)-constituent of \(M_1\oplus M_2\). Since \(M_2\cong M_1^\tau\), both \(M_i\) have
a trivial constituent. Writing \(M_2^{\mathrm{ss}}=\mathbf 1\oplus\alpha\), we get
\(M_1^{\mathrm{ss}}=\bar\eta\oplus\bar\eta\alpha\ni\mathbf 1\), so \(\alpha=\bar\eta^{-1}\). All
four constituents \(\mathbf 1,\bar\eta^{-1},\bar\eta,\mathbf 1\) of \(T'/pT'\) are trivial on
\(G_L\), and \hyperref[lemma-pol]{Lemma~4} applies. \qed

\medskip
We apply these lemmas with \(p=17\), assuming \(A(\Q)[17]\neq0\). Note that the squares in
\(\F_{17}^\times\) are \(1,2,4,8,9,13,15,16\); so \(-1\equiv4^2\) is a square, while
\(-3\equiv14\) and \(-1/3\equiv11\) are not. These residues are also verified in part (B) of the
script.
\begin{itemize}
\item \textbf{\(\Gamma\in\{C_3,C_6,D_6,D_{12}\}\).} Let \(n=3\) if \(\Gamma\in\{C_3,D_6\}\) and
\(n=6\) if \(\Gamma\in\{C_6,D_{12}\}\), and let \(r\in\Gamma\) have order \(n\). The subgroup \(\langle r\rangle\) is
normal, so \(grg^{-1}\in\{r,r^{-1}\}\) for every \(g\in\Gamma\). Let \(R\) be the trace-one
representative of \(r\); its characteristic polynomial is \(t^2-t+1/c\) with \(c=1\) for \(n=3\)
and \(c=3\) for \(n=6\). For \(g\in\Gamma\), the element \(\mathrm{Ad}(R_g)R\) has trace one and
represents \(grg^{-1}\), so it is \(R\) or \(1-R\). Hence \(\mathcal B=\Z_{17}[R]\) is normalized by
\(\Gamma\). The discriminant \(-3\), respectively \(-1/3\), of \(t^2-t+1/c\) is a \(17\)-adic
unit and a non-square modulo \(17\); so \(\mathcal B\) is an unramified quadratic order, \(R\) is a unit
of \(\mathcal B\) (its norm is \(1/c\)), and the reduction of \(R\) lies in \(\F_{289}\setminus\F_{17}\).
\hyperref[lemma-unram]{Lemma~5} with \(g_0=r\) gives \(\Q(\zeta_{17})\subseteq L\). This is
impossible, because \([L:\Q]\le12<16\).

\item \textbf{\(\Gamma=D_8\).} Let \(r\) be a rotation of order \(4\) and \(R\) its trace-one
representative, with characteristic polynomial \(t^2-t+\tfrac12\) of discriminant \(-1\). This
polynomial has two roots \(\lambda_1\equiv11\) and \(\lambda_2\equiv7\pmod{17}\) in \(\Z_{17}\).
So \(\mathcal B=\Z_{17}[R]\), which is normalized by \(\Gamma\) as before, is split, with idempotents
\(e_1=(R-\lambda_2)/(\lambda_1-\lambda_2)\) and \(e_2=(\lambda_1-R)/(\lambda_1-\lambda_2)\). The
rotations fix \(R\), hence \(e_1\) and \(e_2\); the reflections map \(R\) to \(1-R\), and
exchange \(e_1\) and \(e_2\) because \(1-\lambda_2=\lambda_1\). So \(K\) is the fixed field of
\(\langle r\rangle\). For \(\xi\in e_1(D\otimes\Q_{17})e_2\) we have
\(\mathrm{Ad}(R)\xi=(\lambda_1/\lambda_2)\xi\), and \(\lambda_1/\lambda_2\equiv4\pmod{17}\) is
not \(1\). \hyperref[lemma-split]{Lemma~6} gives \(\Q(\zeta_{17})\subseteq L\), which is
impossible because \([L:\Q]=8\).

\item \textbf{\(\Gamma=V_4\).} The map \([M]\mapsto v_{17}(\det M)\bmod2\) is a well-defined
homomorphism \(\Gamma\to\Z/2\); choose \(r\neq1\) in its kernel. A representative \(R\) of \(r\)
has \(R^2=d\in\Q^\times\) and is not a scalar, so \(\tr R=0\) and \(d=-\det R\); thus
\(v_{17}(d)\) is even, and we scale \(R\) so that \(d\) is a \(17\)-adic unit. Since \(\Gamma\) is
abelian, every \(g\in\Gamma\) maps \(R\) to a rational multiple \(\lambda R\) with
\(\lambda^2=1\). So \(\mathcal B=\Z_{17}[R]\) is an \'etale order normalized by \(\Gamma\). If \(d\) is not
a square modulo \(17\), then \(\mathcal B\) is unramified and \(R\) is a unit of \(\mathcal B\) with reduction in
\(\F_{289}\setminus\F_{17}\); \hyperref[lemma-unram]{Lemma~5} applies with \(g_0=r\). If \(d\) is
a square modulo \(17\), then \(d=\delta^2\) with \(\delta\in\Z_{17}^\times\), and \(\mathcal B\) is split with
idempotents \(e_1=(1+R/\delta)/2\) and \(e_2=(1-R/\delta)/2\). Some \(g\in\Gamma\) maps \(R\) to \(-R\):
otherwise \(\Gamma\) would lie in the image of the centralizer \(\Q_{17}[R]^\times\) of \(R\),
that is, in \(\Q_{17}[R]^\times/\Q_{17}^\times\cong\Q_{17}^\times\), whose finite subgroups are
cyclic. This \(g\) exchanges \(e_1\) and \(e_2\), the subgroup fixing them is \(\{1,r\}\), and
\(r\) acts on \(e_1(D\otimes\Q_{17})e_2\) by \(\delta/(-\delta)=-1\). \hyperref[lemma-split]{Lemma~6}
applies. In both cases \(\Q(\zeta_{17})\subseteq L\), which is impossible because
\([L:\Q]=4\).
\end{itemize}

Together with Section~\ref{small-degree} this shows: if \(A\) is of type (S) and
\(A(\Q)[17]\neq0\), then \(\Gamma=C_4\).

\section{The case \(\Gamma=C_4\)}\label{c4}

In this section \(A\) is of type (S), \(\Gamma=\langle r\rangle\cong C_4\), \(s=r^2\), and \(K\)
is the quadratic subfield of \(L\), so that \(\Gal(L/K)=\{1,s\}\). The trace-one representative
\(R\) of \(r\) has \(\det R=\tfrac12\). Put \(J=2R-1\); then \(J^2=-1\), \(R=(1+J)/2\) and
\(R^2=J/2\). So \(s\) acts on \(D\) by \(\mathrm{Ad}(J)\), and \(\End^0_K(A_K)\) is the
centralizer \(\Q(J)\cong\Q(i)\) of \(J\), a field: \textbf{\(A_K\) is simple over \(K\)}.
Moreover \textbf{\(K\) is real}: complex conjugation restricts to an element of order at most
\(2\) of the cyclic group \(\Gamma\), hence to an element of \(\{1,s\}=\Gal(L/K)\).

We assume that \(A(\Q)[17]\neq0\) and derive a contradiction.

\begin{statement}{Lemma 7 (degree class).}\label{lemma-degree}
Let \(X\) be an elliptic curve over \(\bar\Q\) without CM and \(h:X^2\to A_{\bar\Q}\) a
quasi-isogeny. For \(g\in G_{\Q}\), the curve \(X^g\) is isogenous to \(X\); let \(\mu_g\) be a
generator of \(\Hom(X^g,X)\cong\Z\). If \(\det R_g>0\), then
\(\deg\mu_g\equiv\det R_g\pmod{\Q^{\times2}}\).
\end{statement}

\emph{Proof.} The conjugate \(g(h):(X^g)^2\to A_{\bar\Q}\) is a quasi-isogeny, so
\(X^g\sim X\). Put \(N_g=h^{-1}\circ g(h)\). It lies in \(\Hom^0((X^g)^2,X^2)\), whose elements
are the \(N_0\circ\mathrm{diag}(\mu_g,\mu_g)\) with \(N_0\in M_2(\Q)\); so
\(N_g=N_0\circ\mathrm{diag}(\mu_g,\mu_g)\) with \(N_0\in\GL_2(\Q)\). For
\(Y\in M_2(\Q)=\End^0(X^2)\), the conjugate \(g(Y)\in\End^0((X^g)^2)\) is given by the same
matrix, and \(\mathrm{diag}(\mu_g,\mu_g)\circ g(Y)=Y\circ\mathrm{diag}(\mu_g,\mu_g)\). Hence
\[g(hYh^{-1})=g(h)\,g(Y)\,g(h)^{-1}=h\,N_0YN_0^{-1}\,h^{-1}.\]
So, in the identification \(D\cong M_2(\Q)\) given by \(h\), \(g\) acts by \(\mathrm{Ad}(N_0)\).
Therefore \(N_0\) is a rational multiple of the matrix of \(R_g\), and
\(\det N_0\equiv\det R_g\) modulo \(\Q^{\times2}\); in particular \(\det N_0>0\). Since
\(\deg g(h)=\deg h\), (F1) gives
\(1=\deg N_g=(\det N_0)^2(\deg\mu_g)^2\), so \(\deg\mu_g=1/\det N_0\). \qed

\medskip
\textbf{Step 1: a curve over \(K\).}
\begin{enumerate}
\item Let \(E\subset A_L\) be an elliptic subvariety. Then \(E^s\neq E\), since otherwise \(E\)
would be defined over \(K\). So \(A_L\sim_LE\times E^s\). The curve \(E\) has no CM.
\item Apply \hyperref[lemma-degree]{Lemma~7} with \(X=E\) and a lift \(g\) of \(s\): since
\(\det R_g=\det R^2=\tfrac14\), a generator \(\kappa:E^s\to E\) of \(\Hom(E^s,E)\) has square
degree \(m^2\). It is defined over \(L\) (Section~\ref{cases}), and its kernel is cyclic, since
\(\kappa\) is not divisible by an integer greater than \(1\).
\item \emph{The middle curve.} Let \(H_1\subset\ker\kappa\) be the subgroup of order \(m\), which
is \(G_L\)-stable, and \(E_c=E^s/H_1\). Then \(\kappa\) factors as
\(E^s\xrightarrow{\ \kappa_1\ }E_c\xrightarrow{\ \kappa_2\ }E\) through two cyclic isogenies of degree
\(m\). The conjugate \(s(\kappa):E\to E^s\) and the dual \(\hat\kappa\) both generate
\(\Hom(E,E^s)\cong\Z\), so \(s(\kappa)=\pm\hat\kappa\). The unique subgroup of order \(m\) of
\(\ker s(\kappa)=\ker\hat\kappa\) is both \(s(H_1)\) and the kernel of \(\hat\kappa_2:E\to E_c\).
Hence \(E_c^s=E/s(H_1)\cong E_c\), and \(j(E_c)\in K\).
\item Let \(E_K\) be an elliptic curve over \(K\) with \(j(E_K)=j(E_c)\). It has no CM, so
\(j(E_K)\neq0,1728\), and \(A_{\bar\Q}\sim E_K^2\). Fix a quasi-isogeny
\(f:E_K^2\to A_{\bar\Q}\) and let \(\rho_K:G_K\to\GL_2(\Q)\) be the homomorphism of (F3) for
\(k=K\). By (F3)(ii), \(\mathrm{Ad}\circ\rho_K\) is the action of \(G_K\) on \(D\); its kernel is
\(G_L\), so the finite group \(G=\rho_K(G_K)\) has an image of order \(2\) in \(\PGL_2(\Q)\).
\item \emph{\(G\) is cyclic of order \(4\).} The scalars in \(G\) have finite order, so they lie
in \(\{\pm1\}\). Let \(g_0\in G\) be a non-scalar element. By item~4, \(G\) is generated by
\(g_0\) and its scalars, so \(G\) is abelian, and \(g_0^2\) is a scalar: \(g_0^2\in\{\pm1\}\). If
\(g_0^2=1\), the idempotents \((1\pm g_0)/2\) would be nontrivial elements of the commutant of
\(G\), which is isomorphic to the field \(\End^0_K(A_K)\) by (F3)(ii); this is impossible. So
\(g_0^2=-1\) and \(G=\langle g_0\rangle\).
Fix \(\iota\in\Z_{17}\) with \(\iota^2=-1\), and let \(\psi:G_K\to\mu_4(\Z_{17})\) be the
character with \(\psi(g)=\iota^k\) when \(\rho_K(g)=g_0^k\); it has order \(4\). Decomposing
\(\Q_{17}^2\) into the eigenlines of \(g_0\), (F3)(iii) gives an isomorphism of
\(G_K\)-modules
\[V_{17}(A)\cong V_{17}(E_K)\otimes\psi\ \oplus\ V_{17}(E_K)\otimes\psi^{-1}.\]
\end{enumerate}

\textbf{Step 2: a \(2\)-isogeny to the conjugate curve.}
\begin{enumerate}
\item Choose \(\tau\in G_{\Q}\) with \(\tau|_K\neq\mathrm{id}\); then \(\tau|_L\in\{r,r^3\}\)
and \(\tau^2\in G_K\). The curve \(E_K^\tau\) is defined over \(K\), and
\((E_K^\tau)^\tau=E_K\). \hyperref[lemma-degree]{Lemma~7} with \(X=E_K\), \(h=f\) and
\(R_\tau\in\{R,1-R\}\), of determinant \(\tfrac12\), shows that a generator \(\mu:E_K^\tau\to E_K\) of
\(\Hom(E_K^\tau,E_K)\) has degree \(2n^2\) for an integer \(n\ge1\).
\item For \(g\in G_K\), \(g(\mu)\) is again a generator, so \(g(\mu)=\varepsilon(g)\mu\) with a
character \(\varepsilon:G_K\to\{\pm1\}\). So the cyclic group \(\ker\mu\) of order \(2n^2\) and
all its subgroups are \(G_K\)-stable. Let \(Y_1\) and \(Y_2\) be the quotients of \(E_K^\tau\)
by the subgroups of order \(n\) and \(2n\) of \(\ker\mu\); they are elliptic curves over \(K\),
and
\[\mu=c\circ b\circ a,\qquad E_K^\tau\xrightarrow{\ a\ }Y_1\xrightarrow{\ b\ }Y_2
\xrightarrow{\ c\ }E_K,\]
with cyclic isogenies \(a,b,c\) of degrees \(n,2,n\). Here \(a\) and \(b\) are defined over
\(K\), and \(g(c)=\varepsilon(g)c\) for \(g\in G_K\).
\item \emph{A commutative square.} As in Step 1, \(\tau(\mu):E_K\to E_K^\tau\) and
\(\hat\mu\) both generate \(\Hom(E_K,E_K^\tau)\), so \(\tau(\mu)=\pm\hat\mu\). Now
\(\tau(\mu)=\tau(c)\circ\tau(b)\circ\tau(a)\) and \(\hat\mu=\hat a\circ\hat b\circ\hat c\). The
kernels of \(\tau(a)\) and of \(\hat c\) are both the subgroup of order \(n\) of the cyclic group
\(\ker\tau(\mu)=\ker\hat\mu\), and the kernels of \(\tau(b)\circ\tau(a)\) and of
\(\hat b\circ\hat c\) are both its subgroup of order \(2n\). So there are unique isomorphisms
\[\theta_1:Y_1^\tau\to Y_2,\quad\theta_1\circ\tau(a)=\hat c,\qquad\qquad
\theta_2:Y_2^\tau\to Y_1,\quad\theta_2\circ\tau(b)\circ\tau(a)=\hat b\circ\hat c.\]
Since \(\tau(a)\) is surjective, \(\theta_2\circ\tau(b)=\hat b\circ\theta_1\). Taking duals (the
dual of an isomorphism is its inverse) we get
\[\tau(\hat b)\circ\theta_2^{-1}=\theta_1^{-1}\circ b.\]
\item \emph{Notation.} Put \(E_0=Y_2\), an elliptic curve over \(K\) without CM; let
\(\nu=\hat b:E_0\to Y_1\), a \(K\)-rational isogeny of degree \(2\), and
\(\varphi=\theta_2^{-1}:Y_1\to E_0^\tau\), an isomorphism over \(\bar\Q\). The last formula
reads \(\tau(\nu)\circ\varphi=\theta_1^{-1}\circ\hat\nu\), and therefore
\[\varphi(\ker\hat\nu)=\ker\tau(\nu)=\tau(\ker\nu).\]
Both \(Y_1\) and \(E_0^\tau\) are defined over \(K\) and have \(j\neq0,1728\), so
\(g(\varphi)=\eta(g)\varphi\) for \(g\in G_K\), with a character \(\eta:G_K\to\{\pm1\}\).
Finally \(c:E_0\to E_K\) induces an isomorphism \(V_{17}(E_K)\cong V_{17}(E_0)\otimes\varepsilon\)
of \(G_K\)-modules, so that
\[V_{17}(A)\cong V_0\otimes\psi_1\ \oplus\ V_0\otimes\psi_2,\qquad V_0=V_{17}(E_0),\quad
\psi_1=\varepsilon\psi,\quad\psi_2=\varepsilon\psi^{-1},\]
and \(\psi_1,\psi_2\) have order \(4\), because \(\psi_i^2=\psi^{\pm2}\neq1\); note that
\(\psi_2=\psi_1^{-1}\), since \(\varepsilon^2=1\).
\end{enumerate}

\textbf{Step 3: the \(17\)-structure.} Write \(\bar\chi=\bar\chi_{17}|_{G_K}\); it has order
\(16\) or \(8\), because \(\bar\chi_{17}(G_K)\cong\Gal(\Q(\zeta_{17})/K\cap\Q(\zeta_{17}))\) and
\(K\cap\Q(\zeta_{17})\) is \(\Q\) or \(\Q(\sqrt{17})\).
\begin{enumerate}
\item The lattice \(T_{17}(E_0)\otimes\psi_1\oplus T_{17}(E_0)\otimes\psi_2\) has reduction
\(E_0[17]\otimes\bar\psi_1\oplus E_0[17]\otimes\bar\psi_2\). A rational point of order \(17\)
gives a trivial constituent of \(A[17]\), so by (F2) over \(G_K\), \(E_0[17]^{\mathrm{ss}}\)
contains \(\alpha=\bar\psi_1^{-1}\) or \(\alpha=\bar\psi_2^{-1}\), a character of order \(4\)
(reduction is injective on \(\mu_4(\Z_{17})\)). Hence \(E_0[17]\) is reducible over \(K\); let
\(C\subset E_0[17]\) be a \(G_K\)-stable line, with character \(\beta\). By the Weil pairing,
\(E_0[17]^{\mathrm{ss}}=\beta\oplus\bar\chi\beta^{-1}\), and
\(\alpha\in\{\beta,\bar\chi\beta^{-1}\}\).
\item As \(G_K\)-modules, \(E_0^\tau[17]\cong E_0[17]^\tau\) (Section~\ref{facts}). On the other
hand the isogeny \(\varphi\circ\nu:E_0\to E_0^\tau\) has degree \(2\) and satisfies
\(g(\varphi\circ\nu)=\eta(g)\,\varphi\circ\nu\) for \(g\in G_K\), so it induces
\(E_0^\tau[17]\cong E_0[17]\otimes\eta\). Comparing semisimplifications, and using
\(\bar\chi^\tau=\bar\chi\), we get \(\beta^\tau\in\{\beta\eta,\ \bar\chi\beta^{-1}\eta\}\).
\item The second option is impossible: it gives \(\beta\beta^\tau=\bar\chi\eta\). If
\(\beta=\alpha\), the left side has order dividing \(4\); if \(\beta=\bar\chi\alpha^{-1}\), it
gives \(\bar\chi=\alpha\alpha^\tau\eta\). Either way \(\bar\chi^4=1\), whereas \(\bar\chi\) has
order \(16\) or \(8\).
\item Hence \(\beta^\tau=\beta\eta\). Also \(\beta\neq\bar\chi\beta^{-1}\), since
\(\bar\chi=\beta^2\) would give \(\beta=\alpha\) and again \(\bar\chi^4=1\). So the two
constituents of \(E_0^\tau[17]\) are distinct, and \(E_0^\tau[17]\) has at most one
\(G_K\)-stable line with character \(\beta^\tau\). Both \(\tau(C)\) and \(\varphi(\nu(C))\) are
such lines: \(G_K\) acts on the first by \(\beta^\tau\) and on the second by
\(\beta\eta=\beta^\tau\). Therefore
\[\varphi(\nu(C))=\tau(C).\]
\end{enumerate}

\textbf{Step 4: a point of \(X_0(34)\).} We use the description of \(X_0(34)\) and of its
Atkin--Lehner involution \(w_2\) recalled in Section~\ref{x034}, (M2) and (M3). The subgroup
\(C+\ker\nu\subset E_0\) is cyclic of order \(34\) and \(G_K\)-stable, so by (M2)
\[x=(E_0,\ C+\ker\nu)\]
is a non-cuspidal point of \(X_0(34)(K)\). By (M3),
\[w_2(x)=\bigl(E_0/\ker\nu,\ \nu(C)+\nu(E_0[2])\bigr)=\bigl(Y_1,\ \nu(C)+\ker\hat\nu\bigr).\]
Applying the isomorphism \(\varphi:Y_1\to E_0^\tau\) and using Step 3 and Step 2,
\[w_2(x)=\bigl(E_0^\tau,\ \tau(C)+\tau(\ker\nu)\bigr)=x^\tau=x^\sigma,\]
where \(\sigma=\tau|_K\) is the nontrivial automorphism of \(K\). By
\hyperref[prop-x034]{Proposition~9} below, \(K=\Q(\sqrt{-15})\). This contradicts the fact that
\(K\) is real. Hence \(A(\Q)[17]=0\).

\section{The curve \(X_0(34)\)}\label{x034}

\textbf{Standard facts.} Let \(j\) be the modular \(j\)-function on the upper half-plane
\(\mathfrak H\), with \(q\)-expansion \(q^{-1}+744+196884q+\cdots\) in \(q=e^{2\pi iz}\), and put
\(j_d(z)=j(dz)\) for an integer \(d\ge1\). We use the following standard facts about the modular
curve \(X_0(N)\); see for example \cite[Section~1.5 and Sections~7.5--7.7]{DS} and
\cite[Part~II]{DI}.
\begin{itemize}
\item \textbf{(M1)} The field of modular functions for \(\Gamma_0(N)\) is \(\mathbb C(j,j_N)\).
Its degree over \(\mathbb C(j)\) is the index of \(\Gamma_0(N)\) in \(\mathrm{SL}_2(\Z)\), which
is \(54\) for \(N=34\). The curve \(X_0(N)\) is the smooth projective curve over \(\Q\) with
function field \(\Q(j,j_N)\); its complex points are \(\Gamma_0(N)\backslash\mathfrak H^*\), and
its cusps are the poles of \(j\). For a prime \(N\) there are two cusps, \(\infty\) and \(0\);
they are exchanged by the Fricke involution \(z\mapsto-1/(Nz)\), which exchanges \(j\) and
\(j_N\).
\item \textbf{(M2)} The non-cuspidal points of \(X_0(N)(\bar\Q)\) correspond to the
\(\bar\Q\)-isomorphism classes of pairs \((E,C)\), where \(E\) is an elliptic curve over
\(\bar\Q\) and \(C\subset E\) is a cyclic subgroup of order \(N\). The correspondence is
compatible with the action of \(G_{\Q}\), where \((E,C)^g=(E^g,g(C))\); in particular a pair
defined over a number field \(K\) gives a point of \(X_0(N)(K)\). (For pairs without CM, the only
ones used below, this is proved after the list.) At the point \((E,C)\) the
function \(j\) takes the value \(j(E)\) and, for \(d\mid N\), the function \(j_d\) takes the value
\(j(E/C[d])\). Over \(\mathbb C\), the orbit of \(z\in\mathfrak H\) corresponds to the pair
\((\mathbb C/(\Z+\Z z),\langle1/N\rangle)\).
\item \textbf{(M3)} The matrix \(W_2=\begin{pmatrix}2&1\\34&18\end{pmatrix}\), of determinant
\(2\), normalizes \(\Gamma_0(34)\), and \(z\mapsto W_2z\) induces the Atkin--Lehner involution
\(w_2\) of \(X_0(34)\). It is an automorphism of order \(2\) defined over \(\Q\), and on
non-cuspidal points it is given by
\[w_2(E,C)=\bigl(E/C[2],\ (C+E[2])/C[2]\bigr);\]
in particular \(j\circ w_2=j_2\).
\end{itemize}

\emph{The compatibility in (M2) for pairs without CM.} (a) If \(E\) has no CM, then \((E,C)\) is
determined up to isomorphism by \(j(E)\) and \(j(E/C)\). Indeed, let \(C,C'\subset E\) be cyclic
of order \(N\) with \(E/C\cong E/C'\), let \(\pi,\pi'\) be the quotient maps and
\(\theta:E/C\to E/C'\) an isomorphism. Then \(\hat\pi'\circ\theta\circ\pi\in\End(E)=\Z\) has
degree \(N^2\), so it equals \(\pm N=\pm\hat\pi'\circ\pi'\). Hence
\(\hat\pi'\circ(\theta\circ\pi\mp\pi')=0\); as \(\hat\pi'\) has finite kernel and \(E\) is
connected, \(\theta\circ\pi=\pm\pi'\), and \(C=\ker\pi=\ker\pi'=C'\). (b) Since \(j\) and \(j_N\)
lie in \(\Q(X_0(N))\), we have \(h(g(P))=g(h(P))\) for \(h\in\{j,j_N\}\), \(g\in G_{\Q}\) and
\(P\in X_0(N)(\bar\Q)\) not a cusp. If \(P\) corresponds to \((E,C)\) with \(E\) without CM, then
\(g(P)\) and the point corresponding to \((E^g,g(C))\) both have \(j\)-value \(g(j(E))=j(E^g)\)
and \(j_N\)-value \(g(j(E/C))=j(E^g/g(C))\); since \(E^g\) has no CM, they coincide by (a). In
particular a pair without CM defined over \(K\) gives a point fixed by \(G_K\), that is, a point
of \(X_0(N)(K)\).

The two descriptions of \(w_2\) in (M3) agree, as the following computation shows. (That \(w_2\)
is defined over \(\Q\) is not used below; it also follows from
\hyperref[lemma-model]{Lemma~8}.) Put
\(z'=W_2z=(2z+1)/(34z+18)\). Since \((34z+18)-17(2z+1)=1\), we have
\(\Z(34z+18)+\Z(2z+1)=\Z+\Z\cdot2z\), so multiplication by \(34z+18\) is an isomorphism
\(\mathbb C/(\Z+\Z z')\to\mathbb C/(\Z+\Z\cdot2z)\). It maps \(1/34\) to \(z+9/17\), which
generates the subgroup \(\langle z\rangle+\langle1/17\rangle\), because
\(17(z+9/17)\equiv z\) and \(2(z+9/17)\equiv1/17\) modulo \(\Z+\Z\cdot2z\). On the other hand,
for \(E=\mathbb C/(\Z+\Z z)\) and \(C=\langle1/34\rangle\), multiplication by \(2\) is an
isomorphism \(E/C[2]=\mathbb C/(\tfrac12\Z+\Z z)\to\mathbb C/(\Z+\Z\cdot2z)\); it maps
\((C+E[2])/C[2]\), which is generated by \(1/34\) and \(z/2\), onto
\(\langle1/17\rangle+\langle z\rangle\).

\medskip
\textbf{The script.} The computations of this proof are carried out by the script
\href{https://genus-2-torsion.github.io/data/knowledge/sources/geometric_17_checks.m}{\texttt{geometric\_\allowbreak{}17\_\allowbreak{}checks.m}}
(\href{https://genus-2-torsion.github.io/data/knowledge/sources/geometric_17_checks.log}{\texttt{geometric\_\allowbreak{}17\_\allowbreak{}checks.log}}),
which was run in Magma V2.30-1 \cite{Magma} on 10 November 2026. It loads no files and uses only
Magma's built-in databases, it asserts every value quoted in this proof, and it ends with
\texttt{ALL\_CHECKS\_PASSED}. The SHA-256 checksums of the script and of the log are
\texttt{15b9fcd1144ce0f7\allowbreak e2d6677831f1aebd\allowbreak b62b33c416cfb330\allowbreak bcda60c786d21de0}
and
\texttt{a4c7c5a9d8873ebd\allowbreak 795c5a68d33bb93b\allowbreak 2f56700d8e1460f7\allowbreak d0c472d96a5d8bb2}. Its parts (A) and (B) are
the checks quoted in Sections~\ref{good-reduction} and~\ref{lattices}. Part (C) concerns the
plane quartic
\[\mathcal X:\quad x^4-y^4+x^3z+3xy^2z-2x^2z^2+xz^3+z^4=0,\]
which is the model of \(X_0(34)\) in Magma's database of small modular curves, together with the
two maps
\[\omega_2:(x:y:z)\mapsto(z:-y:x),\qquad\omega_{17}:(x:y:z)\mapsto(x:-y:z)\]
and the rational function \(j_{\mathcal X}\) on \(\mathcal X\) that Magma provides as the
\(j\)-function of this model. (In the script these are \texttt{X}, \texttt{w2}, \texttt{w17} and
\texttt{j}.) We do not rely on the correctness of this database: the script verifies the
following facts, and \hyperref[lemma-model]{Lemma~8} deduces from them that
\((\mathcal X,j_{\mathcal X},\omega_2)\) is a model of \((X_0(34),j,w_2)\).
\begin{itemize}
\item \textbf{(C0)} For \(N\in\{2,17\}\), let \(\Phi_N(u,v)\in\Z[u,v]\) be the polynomial
returned by Magma's \texttt{ClassicalModularPolynomial}. It has degree \(N+1\) in each variable
and is monic in \(v\). The script computes the \(q\)-expansion of \(j=E_4^3/\Delta\) from the
definitions of \(E_4\) and \(\Delta\), and finds \(\Phi_2(j,j_2)=O(q^{72})\) and
\(\Phi_{17}(j,j_{17})=O(q^{402})\).
\item \textbf{(C1)} \(\mathcal X\) is a nonsingular plane quartic; so it is a smooth projective
geometrically irreducible curve of genus \(3\).
\item \textbf{(C2a)} \(\omega_2\) and \(\omega_{17}\) are involutions of \(\mathcal X\), the
function \(j_{\mathcal X}\) has degree \(54\), and
\(\Phi_2(j_{\mathcal X},j_{\mathcal X}\circ\omega_2)=0\); the last identity is an exact
computation in the function field of \(\mathcal X\). The degree is computed twice: by Magma's
function field arithmetic, and as \(4\cdot53-158\): here \(j_{\mathcal X}=j_{\rm num}/j_{\rm den}\)
with forms \(j_{\rm num},j_{\rm den}\) of degree \(53\), and \(158\) is the length of the base scheme
\(\mathcal X\cap\{j_{\rm num}=j_{\rm den}=0\}\). Thus \(j_{\mathcal X}\) has \(54\) poles counted with multiplicity.
\item \textbf{(C2c)} For each integer \(2\le x_0\le502\), the polynomial \(f(x_0,Y,1)\), where
\(f\) is the equation of \(\mathcal X\), is a squarefree quartic. At each of the resulting
\(4\cdot501=2004\) geometric points \((x_0:y_0:1)\) of \(\mathcal X\), the functions
\(j_{\mathcal X}\) and \(j_{\mathcal X}\circ\omega_{17}\) are regular, and
\(\Phi_{17}(j_{\mathcal X},j_{\mathcal X}\circ\omega_{17})\) vanishes. (The script works in the
number fields generated by the roots \(y_0\).)
\item \textbf{(C3)} The quotient of \(\mathcal X\) by \(\omega_2\) is a nonsingular curve
\(\mathcal Y\subset\mathbb P^3\) of genus \(1\), with a morphism \(\pi:\mathcal X\to\mathcal Y\)
of degree \(2\) defined over \(\Q\) such that \(\pi\circ\omega_2=\pi\). The curve \(\mathcal Y\)
is isomorphic over \(\Q\) to the elliptic curve
\[\text{17a2}:\quad y^2+xy+y=x^3-x^2-6x-4,\]
whose torsion subgroup is \(\Z/2\times\Z/2\) and whose \(2\)-Selmer group has order \(4\). So
its rank is \(0\), and \(\mathcal Y\) has exactly four rational points. This is unconditional:
the curve has full rational \(2\)-torsion, so no class group of a number field enters the
\(2\)-descent.
\item \textbf{(C4)} The fibres of \(\pi\) over the four rational points of \(\mathcal Y\) are
as follows. Two fibres consist of two rational points each; these four points are
\[(0:1:1),\quad(-1:1:0),\quad(0:-1:1),\quad(1:1:0),\]
and \(j_{\mathcal X}\) has poles of order \(34,17,2,1\) at them, respectively. Each of the
other two fibres is a single closed point of degree \(2\) with residue field
\(\Q(\sqrt{-15})\), namely
\[\Bigl(\tfrac{7\pm\sqrt{-15}}8:\tfrac{1\mp\sqrt{-15}}8:1\Bigr)\qquad\text{and}\qquad
\Bigl(\tfrac{7\pm\sqrt{-15}}8:\tfrac{-1\pm\sqrt{-15}}8:1\Bigr).\]
At these points \(j_{\mathcal X}\) is finite, and its values have the minimal polynomials
\[Y^2-\tfrac{11779}{4}\,Y+\tfrac{25153757}{8}\qquad\text{and}\qquad
Y^2+\tfrac{7319387769191}{17179869184}\,Y+\tfrac{81029690385834182553917}{2251799813685248},\]
one for each of the two closed points.
\item \textbf{(C4\('\))} The involution \(\omega_2\) maps each of the four points over
\(\Q(\sqrt{-15})\) in (C4) to its conjugate. (This also follows from (C3) and (C4): \(\omega_2\)
preserves the fibres of \(\pi\) and has no fixed point in a fibre with two geometric points.)
\end{itemize}
The script contains two further consistency checks: in (C1) it compares the number of points of
\(\mathcal X\) over \(\F_p\) with the traces of the Hecke operators on \(S_2(\Gamma_0(34))\) for
the primes \(3\le p<100\), \(p\neq17\), and in (C2b) it tests \(\omega_2\) and \(\omega_{17}\)
against \(\Phi_2\) and \(\Phi_{17}\) at five points over number fields. These two checks and
(C4\('\)) are not used in the proof of the Theorem.

\begin{statement}{Lemma 8 (the model).}\label{lemma-model}
There is an isomorphism \(\Theta:\mathcal X\to X_0(34)\) of curves over \(\Q\) such that
\(j\circ\Theta=j_{\mathcal X}\) and \(\Theta\circ\omega_2=w_2\circ\Theta\). In particular
\(\Theta\) maps the poles of \(j_{\mathcal X}\) onto the cusps of \(X_0(34)\).
\end{statement}

\emph{Proof.}
\begin{enumerate}
\item \emph{\(\Phi_N(j,j_N)=0\) for \(N\in\{2,17\}\).} The function \(g=\Phi_N(j,j_N)\) is a
modular function for \(\Gamma_0(N)\) that is holomorphic on \(\mathfrak H\). In terms of the
local parameter \(q\) at the cusp \(\infty\), \(j\) has a simple pole and \(j_N\) a pole of order
\(N\); by the Fricke involution in (M1), \(j\) has a pole of order \(N\) and \(j_N\) a simple pole at
the cusp \(0\). So a
monomial \(j^aj_N^b\) with \(a,b\le N+1\) has a pole of order at most \((N+1)^2\) at each cusp.
If \(g\) were nonzero and vanished at \(\infty\), all its poles would be at the cusp \(0\), and
\(g\) would have at most \((N+1)^2\) zeros, counted with multiplicity. But by (C0), \(g\)
vanishes at \(\infty\) to order at least \(72>9\) for \(N=2\) and at least \(402>324\) for
\(N=17\). Hence \(g=0\).

\item \emph{\(\Q(X_0(34))=\Q(j,j_2,j_{17})\), of degree \(54\) over \(\Q(j)\).} For
\(\gamma\in\mathrm{SL}_2(\Z)\) we have \(j_d\circ\gamma=j_d\) if and only if
\(\gamma\in\Gamma_0(d)\): with \(\delta=\mathrm{diag}(d,1)\) we have
\(j_d\circ\gamma=j\circ(\delta\gamma\delta^{-1})\circ\delta\), and this equals \(j\circ\delta\)
exactly when \(\delta\gamma\delta^{-1}\in\mathrm{SL}_2(\Z)\). (If \(j\circ M=j\) for some
\(M\in\mathrm{SL}_2(\Q)\), then \(\mathfrak H\) is the union of the countably many closed sets
\(\{z:\gamma^{-1}Mz=z\}\), \(\gamma\in\mathrm{SL}_2(\Z)\), so some \(\gamma^{-1}M\) fixes a
nonempty open set and equals \(\pm1\); hence \(M=\pm\gamma\in\mathrm{SL}_2(\Z)\).) Hence the pairs
\((j_2\circ\gamma,j_{17}\circ\gamma)\), for \(\gamma\) in a set of representatives of
\(\Gamma_0(34)\backslash\mathrm{SL}_2(\Z)\), are \(54\) distinct pairs, and for all but finitely
many \(c\in\Q\) the \(54\) functions \(\vartheta\circ\gamma\), where
\(\vartheta=j_2+c\,j_{17}\), are distinct. They are roots of the minimal polynomial of
\(\vartheta\) over \(\mathbb C(j)\), whose coefficients are invariant under
\(\mathrm{SL}_2(\Z)\). So \([\mathbb C(j,j_2,j_{17}):\mathbb C(j)]\ge54\); as this field is
contained in \(\mathbb C(j,j_{34})\), of degree \(54\) by (M1), the two fields are equal.

Elements of \(\Q((q))\) that are linearly independent over \(\Q(j)\) remain linearly
independent over \(\mathbb C(j)\): in a relation with coefficients in \(\mathbb C[j]\), expand
the complex coefficients in a basis of a finite-dimensional \(\Q\)-vector space containing
them. So \([M:\Q(j)]=[\mathbb C(j)M:\mathbb C(j)]\) for every field
\(\Q(j)\subseteq M\subset\Q((q))\) that is finite over \(\Q(j)\). The functions
\(j_2,j_{17},j_{34}\) are algebraic over \(\Q(j)\): the first two by step~1, and \(j_{34}\) is a
root of \(\Phi_{17}(j_2,v)\). Therefore the fields \(\Q(j,j_{34})\), \(\Q(j,j_2,j_{17})\) and
\(\Q(j,j_2,j_{17},j_{34})\) all have degree \(54\) over \(\Q(j)\), and they coincide.

\item \emph{The isomorphism \(\Theta\).} The \(\Q(j)\)-algebra
\[\mathcal A=\Q(j)[v_2,v_{17}]/\bigl(\Phi_2(j,v_2),\ \Phi_{17}(j,v_{17})\bigr)\]
has dimension \(3\cdot18=54\) by (C0), and by step~1 it maps onto \(\Q(j,j_2,j_{17})\) via
\(v_d\mapsto j_d\). By step~2 this map is an isomorphism, so \(\mathcal A\) is a field,
isomorphic to \(\Q(X_0(34))\). By (C2a) and (C2c), the assignment
\[j\mapsto j_{\mathcal X},\qquad v_2\mapsto j_{\mathcal X}\circ\omega_2,\qquad
v_{17}\mapsto j_{\mathcal X}\circ\omega_{17}\]
defines a homomorphism from \(\mathcal A\) to the function field \(\Q(\mathcal X)\). For (C2c),
note that \(h=\Phi_{17}(j_{\mathcal X},j_{\mathcal X}\circ\omega_{17})\) has a divisor of poles
of degree at most \(18\cdot54+18\cdot54=1944\), and that it vanishes at \(2004\) points; so
\(h=0\). The homomorphism is injective, and
\([\Q(\mathcal X):\Q(j_{\mathcal X})]=54=[\mathcal A:\Q(j)]\), so it is an isomorphism
\(\Q(X_0(34))\to\Q(\mathcal X)\). It corresponds to an isomorphism
\(\Theta:\mathcal X\to X_0(34)\) over \(\Q\) with \(j\circ\Theta=j_{\mathcal X}\) and
\(j_2\circ\Theta=j_{\mathcal X}\circ\omega_2\).

\item \emph{\(\Theta\circ\omega_2=w_2\circ\Theta\).} Put
\(w=\Theta\circ\omega_2\circ\Theta^{-1}\). Then \(j\circ w=j_2=j\circ w_2\) by step~3 and (M3),
so the automorphism \(u=w\circ w_2^{-1}\) of \(X_0(34)\) satisfies \(j\circ u=j\). Thus \(u\)
permutes the poles of \(j\) and preserves their orders. By (C4), transported by \(\Theta\),
these orders are \(34,17,2,1\); their sum is \(54=\deg j\), so these are all the poles. Hence
\(u\) fixes the simple pole \(P_\infty\) of \(j\). The function \(1/j\) is a uniformizer at
\(P_\infty\) and is fixed by \(u\); so \(u\) acts trivially on the completed local ring at
\(P_\infty\), hence on the function field, and \(u=1\). \qed
\end{enumerate}

\begin{statement}{Proposition 9.}\label{prop-x034}
Let \(K\) be a quadratic field with nontrivial automorphism \(\sigma\), and let
\(x\in X_0(34)(K)\) be a non-cuspidal point with \(w_2(x)=x^\sigma\). Then
\(K=\Q(\sqrt{-15})\).
\end{statement}

\emph{Proof.} Let \(x'=\Theta^{-1}(x)\in\mathcal X(K)\), with \(\Theta\) as in
\hyperref[lemma-model]{Lemma~8}; it is not a pole of \(j_{\mathcal X}\), and
\(\omega_2(x')=x'^\sigma\). Then
\(\pi(x')^\sigma=\pi(x'^\sigma)=\pi(\omega_2(x'))=\pi(x')\), so \(\pi(x')\) is one of the four
rational points of \(\mathcal Y\) of (C3). By (C4), the fibres of \(\pi\) over two of them
consist of poles of \(j_{\mathcal X}\). So \(x'\) lies in one of the other two fibres, each of
which is a single closed point with residue field \(\Q(\sqrt{-15})\). Hence
\(\Q(\sqrt{-15})\subseteq K\), and \(K=\Q(\sqrt{-15})\). \qed

\medskip
Conversely, by (C4\('\)) the four points over \(\Q(\sqrt{-15})\) do satisfy
\(w_2(x)=x^\sigma\). So the fact that \(K\) is real is essential in Section~\ref{c4}. As a by-product, the rational
points of \(X_0(34)\) are its four cusps.

\medskip
\emph{A check by hand of (C3) and (C4).} Put \(a=x+z\) and \(b=x-z\). Then \(\omega_2\) becomes
\((a:b:y)\mapsto(-a:b:y)\), and the equation of \(\mathcal X\) becomes
\[8f=a^4+(8b^2+6y^2)\,a^2-(b^2+2y^2)(b^2+4y^2),\]
a quadratic polynomial in \(a^2\) with discriminant \(4W^2\), where
\(W^2=17b^4+30b^2y^2+17y^4\). The map \((a:b:y)\mapsto(b:y:a^2+4b^2+3y^2)\) identifies the quotient
\(\mathcal X/\omega_2\) with this genus-\(1\) curve. It has the four rational points
\((b:y:W)=(\pm1:1:\pm8)\), and by (C3) no others. Above a rational point,
\(a^2=W-4b^2-3y^2\), which is \(1\) for \(W=8\) and \(-15\) for \(W=-8\). The value \(1\) gives the
four rational points of (C4), and the value \(-15\) gives the points over \(\Q(\sqrt{-15})\).

\section{Proof of the Theorem and consequence for the census}\label{census}

\emph{Proof of the Theorem.} Let \(A/\Q\) be a geometrically split abelian surface. By
Section~\ref{cases} it is of type (N), (CM) or (S). In types (N) and (CM), \(A(\Q)\) has no point
of order \(17\) by Section~\ref{small-degree} and Section~\ref{good-reduction}. In type (S) the
group \(\Gamma\) is one of the nine groups of Section~\ref{cases}: the groups \(C_1\) and \(C_2\)
are treated in Section~\ref{small-degree}, the groups \(C_3,C_6,V_4,D_6,D_8,D_{12}\) in
Section~\ref{lattices}, and \(C_4\) in Section~\ref{c4}, using
\hyperref[prop-x034]{Proposition~9}. \qed

\medskip
\emph{Proof of the Corollary.} The Jacobian of a curve of genus \(2\) over \(\Q\) is an abelian
surface over \(\Q\). \qed

\medskip
The group \href{https://genus-2-torsion.github.io/group.html?g=17}{\([17]\)} is impossible in the
class ``geometrically split, simple over \(\Q\)''. It was already impossible in the class ``split
over \(\Q\)'' by Mazur's theorem
(\href{https://genus-2-torsion.github.io/proofs/qsplit-prime-obstruction.html}{prime torsion
obstruction}). Hence no geometrically split genus-\(2\) Jacobian over \(\Q\) has a rational point
of order \(17\), and the combined geometrically split cell of
\href{https://genus-2-torsion.github.io/group.html?g=17}{\([17]\)}, whose grade was open
(``\(\infty\)?'' on the census), becomes impossible (\(\emptyset\)).

The Theorem applies to every group containing a point of order \(17\). Among the census groups
these are \href{https://genus-2-torsion.github.io/group.html?g=17}{\([17]\)} and
\href{https://genus-2-torsion.github.io/group.html?g=34}{\([34]\)}. The latter was already
excluded in both geometrically split classes by the
\href{https://genus-2-torsion.github.io/proofs/geometric-33-34-obstruction.html}{proof for orders
33 and 34}.

The argument places no restriction on the geometric splitting degree, the field of definition of
the geometric endomorphisms, or the integral endomorphism ring, and it does not use a principal
polarization.

\section{Remarks}\label{remarks}

\begin{enumerate}
\item \emph{Fields of degree \(4\).} By Derickx, Kamienny, Stein and Stoll
\cite[Theorem~1.2]{DKSS}, \(S(4)=\mathrm{Primes}(17)\): there are elliptic curves over quartic
fields with a point of order \(17\). This is why the argument of Section~\ref{small-degree}
cannot be used when \([L:\Q]\ge4\). For \(\Gamma=C_3\), where \([L:\Q]=3\), it could be used,
since \(S(3)=\mathrm{Primes}(13)\) by Parent's theorem, as quoted in \cite[Section~1]{DKSS}; we
treat \(C_3\) together with \(C_6\) in Section~\ref{lattices} instead.

\item \emph{Quadratic points of \(X_0(34)\).} Ozman and Siksek \cite{OS} determined all
quadratic points of \(X_0(34)\) (their Main Theorem and Table~8.1). The non-cuspidal ones are
defined over \(\Q(\sqrt{-1})\), \(\Q(\sqrt{-2})\) and \(\Q(\sqrt{-15})\). Those over
\(\Q(\sqrt{-15})\) are the points \(P_5\), \(P_6\) of their table and the conjugates; they have
no CM, the involution \(w_2\) maps each of them to its conjugate, and their \(j\)-invariants
\((2041\theta+11779)/8\) and \((-53184785340479\,\theta-7319387769191)/34359738368\), with
\(\theta^2=-15\), have the minimal polynomials of (C4). Since none of the three fields is real,
their theorem gives a second proof of the last step of Section~\ref{c4}. The proof above does
not depend on it.

\item \emph{Potentially good reduction.} By Clark \cite{Cla}, the order of the rational torsion
subgroup of an abelian surface over \(\Q\) with potentially good reduction everywhere belongs to
\(\{1,\ldots,16,18,19,20,22,24,25,28,30,36,48,60,72\}\), which contains no multiple of \(17\).
\hyperref[prop-good]{Proposition~3} is the part of this statement that is needed here; it only
uses the reduction at \(2\).

\item \emph{Relation to the earlier proofs.} The Theorem supersedes the last sentence of the
\href{https://genus-2-torsion.github.io/proofs/geometric-33-34-obstruction.html}{proof for orders
33 and 34}, which left rational \(17\)-torsion open. The lemmas of Section~\ref{lattices} and the
construction of Section~\ref{c4} are the analogues for \(p=17\) of those in the
\href{https://genus-2-torsion.github.io/proofs/qsimple-13-obstruction.html}{proof for order 13};
they are given here in full. For \(p=17\) the non-residues \(-3\) and \(-1/3\) dispose of the
groups \(C_3\) and \(C_6\) together with the dihedral groups, so that \(C_4\) is the only group
which requires a modular curve; and, in contrast to order \(13\), neither a principal
polarization nor simplicity over \(\Q\) is needed. The construction of the curve \(E_c\) in
Section~\ref{c4} is a special case of Elkies's reduction of quadratic \(K\)-curves
(Elkies~\cite{Elk}).
\end{enumerate}

\begin{thebibliography}{DKSS}

\bibitem[Cla]{Cla}
P. L. Clark, \emph{Bounds for torsion on abelian varieties with integral moduli},
arXiv:math/0407264v2 (2004), Theorem 5. \url{https://arxiv.org/pdf/math/0407264}.
The same list is Theorem~1.4 of P. L. Clark and X. Xarles, \emph{Local bounds for torsion points
on abelian varieties}, Canad. J. Math. \textbf{60} (2008), 532--555,
\href{https://doi.org/10.4153/CJM-2008-026-x}{doi:10.4153/CJM-2008-026-x}, for abelian surfaces
with everywhere anisotropic reduction, which includes potentially good reduction.

\bibitem[DKSS]{DKSS}
M. Derickx, S. Kamienny, W. Stein and M. Stoll, \emph{Torsion points on elliptic curves over
number fields of small degree}, Algebra Number Theory \textbf{17} (2023), 267--308;
\href{https://arxiv.org/abs/1707.00364}{arXiv:1707.00364}.

\bibitem[DI]{DI}
F. Diamond and J. Im, \emph{Modular forms and modular curves}, in: \emph{Seminar on Fermat's
Last Theorem (Toronto, 1993--1994)}, CMS Conf. Proc. \textbf{17}, Amer. Math. Soc., Providence,
1995, 39--133.

\bibitem[DS]{DS}
F. Diamond and J. Shurman, \emph{A First Course in Modular Forms}, Graduate Texts in Mathematics
\textbf{228}, Springer, New York, 2005.

\bibitem[Elk]{Elk}
N. D. Elkies, \emph{On elliptic \(K\)-curves}, in: \emph{Modular curves and abelian varieties},
Progr. Math. \textbf{224}, Birkh\"auser, Basel, 2004, 81--91.

\bibitem[Kam]{Kam}
S. Kamienny, \emph{Torsion points on elliptic curves and \(q\)-coefficients of modular forms},
Invent. Math. \textbf{109} (1992), 221--229, Theorem~3.1.
\href{https://doi.org/10.1007/BF01232025}{doi:10.1007/BF01232025}.

\bibitem[Magma]{Magma}
W. Bosma, J. Cannon and C. Playoust, \emph{The Magma algebra system I: the user language},
J. Symbolic Comput. \textbf{24} (1997), 235--265.
\href{https://doi.org/10.1006/jsco.1996.0125}{doi:10.1006/jsco.1996.0125}.

\bibitem[Maz]{Maz}
B. Mazur, \emph{Modular curves and the Eisenstein ideal}, Publ. Math. IH\'ES \textbf{47} (1977),
33--186, Theorem~(8). \href{https://www.numdam.org/item/PMIHES_1977__47__33_0/}{Numdam}.

\bibitem[OS]{OS}
E. Ozman and S. Siksek, \emph{Quadratic points on modular curves}, Math. Comp. \textbf{88}
(2019), 2461--2484; \href{https://arxiv.org/abs/1806.08192}{arXiv:1806.08192}.

\bibitem[ST]{ST}
J.-P. Serre and J. Tate, \emph{Good reduction of abelian varieties}, Ann. of Math. (2)
\textbf{88} (1968), 492--517: Theorem~1 (p.~493; the criterion of N\'eron, Ogg and Shafarevich)
and \S2 (p.~496).
\href{https://doi.org/10.2307/1970722}{doi:10.2307/1970722}.

\end{thebibliography}

\end{document}
